% ============================================================
% CAPÍTULO 5: APÉNDICES
% ============================================================
\section{Apéndices}
\label{sec:apendices}

Este capítulo contiene información complementaria que apoya la comprensión del sistema AnxioTest, incluyendo un glosario detallado, formatos de datos, referencias adicionales y documentación de soporte.

\subsection{Glosario}
\label{sec:glosario}

\begin{longtable}{p{0.25\linewidth} p{0.70\linewidth}}
\toprule
\rowcolor{tableheader!20}
\textbf{Término} & \textbf{Definición} \\
\midrule
\endfirsthead
\rowcolor{tableheader!20}
\textbf{Término} & \textbf{Definición} \\
\midrule
\endhead
\bottomrule
\endfoot
\textbf{AnxioTest}
  & Nombre de la plataforma web de evaluación de bienestar emocional desarrollada para VincuMedicina 2026. Combina ``anxiety'' (ansiedad) y ``test'' (prueba), reflejando su propósito de evaluación. \\
\textbf{BAI}
  & Beck Anxiety Inventory. Inventario de Ansiedad de Beck, instrumento de 21 ítems desarrollado por Aaron Beck y colaboradores (1988) para medir la intensidad de síntomas de ansiedad en la población general. \\
\textbf{PHQ-9}
  & Patient Health Questionnaire -- 9 items. Cuestionario de Salud del Paciente de 9 preguntas desarrollado por Kroenke, Spitzer y Williams (2001) para detectar síntomas depresivos en atención primaria. \\
\textbf{Tamizaje}
  & Proceso de identificación preliminar de posibles casos de trastornos mentales mediante instrumentos estandarizados, con el fin de derivar a evaluación profesional cuando sea necesario. \\
\textbf{Síntoma depresivo}
  & Manifestación de alteración del estado de ánimo, incluyendo tristeza, pérdida de interés, cambios en el apetito, alteraciones del sueño, fatiga, sentimientos de culpa o inutilidad, dificultad de concentración y pensamientos de muerte. \\
\textbf{Síntoma de ansiedad}
  & Manifestación de ansiedad, incluyendo síntomas somáticos (palpitaciones, sudoración, temblores) y cognitivos (preocupación excesiva, miedo, sensación de peligro). \\
\textbf{TTS}
  & Text-to-Speech. Tecnología que convierte texto en voz sintetizada utilizando la API Web Speech del navegador. \\
\textbf{ARIA}
  & Accessible Rich Internet Applications. Conjunto de atributos HTML definidos por el W3C que mejoran la accesibilidad para usuarios de lectores de pantalla y otras tecnologías asistivas. \\
\textbf{sessionStorage}
  & Almacenamiento temporal de datos en el navegador durante la sesión actual. Los datos se eliminan al cerrar la pestaña o el navegador. Capacidad típica: 5--10 MB. \\
\textbf{localStorage}
  & Almacenamiento persistente de datos en el navegador que permanece entre sesiones del usuario. Capacidad típica: 5--10 MB. \\
\textbf{SPA}
  & Single Page Application. Arquitectura web que carga una única página HTML y actualiza dinámicamente el contenido sin recarga completa del navegador. \\
\textbf{VincuMedicina}
  & Iniciativa de vinculación con la colectividad de la Universidad de las Fuerzas Armadas ESPE, que promueve proyectos de salud y bienestar comunitario. \\
\textbf{ESPE}
  & Universidad de las Fuerzas Armadas, institución de educación superior del Ecuador que promueve la investigación y vinculación con la sociedad. \\
\textbf{Severidad}
  & Grado de intensidad o gravedad de los síntomas reportados, clasificado en categorías que van desde mínimo hasta grave según los criterios de cada instrumento. \\
\bottomrule
\caption{Glosario de términos}
\label{tab:glosario}
\end{longtable}

\subsection{Formato de datos en sessionStorage}
\label{sec:formato_datos}

El sistema guarda los resultados de evaluación temporales con la clave \texttt{lastTestResult} en formato JSON con la siguiente estructura:

\begin{lstlisting}[caption={Estructura de datos en sessionStorage},
                   label=lst:sessionstorage]
{
  "testId":          "BAI" | "PHQ9",
  "testName":        "BAI - Ansiedad" | "PHQ-9 - Bienestar",
  "totalScore":      15,
  "maxScore":        27,
  "level":           "Categoria de resultado",
  "levelLabel":      "Etiqueta amigable para el usuario",
  "message":         "Mensaje interpretativo personalizado",
  "recommendations": [
    "Recomendacion 1",
    "Recomendacion 2",
    "Recomendacion 3"
  ],
  "answeredCount":   9,
  "totalQuestions":  9,
  "timestamp":       "2025-05-15T10:30:00.000Z"
}
\end{lstlisting}

\begin{longtable}{p{0.20\linewidth} p{0.20\linewidth} p{0.55\linewidth}}
\toprule
\rowcolor{tableheader!20}
\textbf{Campo} & \textbf{Tipo} & \textbf{Descripción} \\
\midrule
\endfirsthead
\rowcolor{tableheader!20}
\textbf{Campo} & \textbf{Tipo} & \textbf{Descripción} \\
\midrule
\endhead
\bottomrule
\endfoot
\texttt{testId}          & String          & Identificador único de la prueba: ``BAI'' o ``PHQ9''. \\
\texttt{testName}        & String          & Nombre completo de la prueba para visualización. \\
\texttt{totalScore}      & Number          & Puntaje total obtenido en el cuestionario. \\
\texttt{maxScore}        & Number          & Puntaje máximo posible para la prueba. \\
\texttt{level}           & String          & Categoría de severidad (ej.\ ``Moderado''). \\
\texttt{levelLabel}      & String          & Etiqueta amigable para mostrar al usuario (ej.\ ``Algunos días pesan más que otros''). \\
\texttt{message}         & String          & Mensaje interpretativo personalizado. \\
\texttt{recommendations} & Array of String & Lista de recomendaciones prácticas. \\
\texttt{answeredCount}   & Number          & Número de preguntas respondidas. \\
\texttt{totalQuestions}  & Number          & Número total de preguntas del cuestionario. \\
\texttt{timestamp}       & String          & Fecha y hora de finalización en formato ISO 8601. \\
\bottomrule
\caption{Campos de datos en sessionStorage}
\label{tab:sessionstorage}
\end{longtable}

\subsection{Documentación de soporte}
\label{sec:documentacion_soporte}

\subsubsection{Actas de reuniones}

Documento que registra las reuniones de seguimiento del proyecto, incluyendo el acta de la primera, segunda y tercera reunión. En estas actas se detallan los acuerdos, decisiones y avances relacionados con la planificación, revisión de requisitos y validación con los stakeholders del proyecto.

\subsubsection{Justificaciones técnicas}

Conjunto de documentos que sustentan las decisiones técnicas adoptadas durante el desarrollo del sistema:
\begin{itemize}
  \item \textbf{Justificación de colores:} Documento que explica la selección de la paleta cromática utilizada en la interfaz, considerando principios de accesibilidad y experiencia de usuario.
  \item \textbf{Justificación breve de resultados y feedback de doctores:} Resumen ejecutivo de las validaciones realizadas con profesionales de la salud.
  \item \textbf{Justificación técnica de resultados y feedback de pruebas:} Documento detallado que analiza los resultados obtenidos en las pruebas funcionales y de usabilidad.
\end{itemize}

\subsubsection{Manuales}

Documento informativo del cuestionario de bienestar emocional, dirigido a los usuarios finales (adultos mayores y comunidad). Este manual explica el propósito de la herramienta, la forma de uso de la plataforma y la interpretación de los resultados obtenidos.

\subsubsection{Manual de pruebas}

Documento en proceso de elaboración que describirá los casos de prueba, los procedimientos de validación y los criterios de aceptación del sistema AnxioTest. Una vez finalizado, permitirá al equipo de pruebas ejecutar y documentar las validaciones funcionales y no funcionales del sistema.

\subsection{Resumen de referencia rápida}
\label{sec:referencia_rapida}

\subsubsection{PHQ-9 --- Puntos de corte}

\begin{longtable}{p{0.15\linewidth} p{0.35\linewidth} p{0.45\linewidth}}
\toprule
\rowcolor{tableheader!20}
\textbf{Puntaje} & \textbf{Nivel clínico estándar} & \textbf{Etiqueta mostrada al paciente} \\
\midrule
\endfirsthead
\rowcolor{tableheader!20}
\textbf{Puntaje} & \textbf{Nivel clínico estándar} & \textbf{Etiqueta mostrada al paciente} \\
\midrule
\endhead
\bottomrule
\endfoot
0 -- 4   & Mínimo / Sin síntomas  & Buen equilibrio de ánimo \\
5 -- 9   & Leve                   & Días con altibajos \\
10 -- 14 & Moderado               & Algunos días pesan más que otros \\
15 -- 19 & Moderadamente grave    & Un poco de compañía podría hacerle bien \\
20 -- 27 & Grave                  & Es tiempo de tomar un respiro \\
\bottomrule
\caption{PHQ-9 --- Puntos de corte}
\label{tab:phq9_corte}
\end{longtable}

\subsubsection{BAI --- Puntos de corte}

\begin{longtable}{p{0.15\linewidth} p{0.35\linewidth} p{0.45\linewidth}}
\toprule
\rowcolor{tableheader!20}
\textbf{Puntaje} & \textbf{Nivel clínico estándar} & \textbf{Etiqueta mostrada al paciente} \\
\midrule
\endfirsthead
\rowcolor{tableheader!20}
\textbf{Puntaje} & \textbf{Nivel clínico estándar} & \textbf{Etiqueta mostrada al paciente} \\
\midrule
\endhead
\bottomrule
\endfoot
0 -- 21  & Ansiedad mínima   & Días de mayor calma \\
22 -- 35 & Ansiedad moderada & Algunos momentos de inquietud \\
36 -- 63 & Ansiedad severa   & Días emocionalmente difíciles \\
\bottomrule
\caption{BAI --- Puntos de corte}
\label{tab:bai_corte}
\end{longtable}

\subsubsection{Recomendaciones por nivel --- PHQ-9}

\begin{longtable}{p{0.12\linewidth} p{0.83\linewidth}}
\toprule
\rowcolor{tableheader!20}
\textbf{Nivel} & \textbf{Recomendaciones} \\
\midrule
\endfirsthead
\rowcolor{tableheader!20}
\textbf{Nivel} & \textbf{Recomendaciones} \\
\midrule
\endhead
\bottomrule
\endfoot
0--4   & Continuar con actividades placenteras, mantener hábitos de cuidado personal, dedicar tiempo a seres queridos, inspirar a otros con el bienestar. \\
5--9   & Dedicar tiempo a actividades placenteras, compartir sentimientos con alguien cercano, mantener horarios regulares, ofrecer palabras amables. \\
10--14 & Hablar con familiar o amigo de confianza, visitar al médico o centro de salud, mantener pequeñas rutinas diarias, practicar respiración profunda. \\
15--19 & Compartir sentimientos con alguien de confianza, buscar orientación profesional, realizar ejercicios de relajación, mantener contacto social. \\
20--27 & Estar cerca de personas de confianza, buscar apoyo profesional comprensivo, recordar que importa a quienes le rodean, buscar compañía. \\
\bottomrule
\caption{Recomendaciones por nivel --- PHQ-9}
\label{tab:phq9_recomendaciones}
\end{longtable}

\subsubsection{Recomendaciones por nivel --- BAI}

\begin{longtable}{p{0.12\linewidth} p{0.83\linewidth}}
\toprule
\rowcolor{tableheader!20}
\textbf{Nivel} & \textbf{Recomendaciones} \\
\midrule
\endfirsthead
\rowcolor{tableheader!20}
\textbf{Nivel} & \textbf{Recomendaciones} \\
\midrule
\endhead
\bottomrule
\endfoot
0--21  & Mantener actividades que generen calma, practicar técnicas de respiración, cultivar relaciones de apoyo y mantener hábitos de sueño regulares. \\
22--35 & Identificar y reducir fuentes de estrés, practicar meditación o relajación muscular progresiva, hablar con alguien de confianza, considerar apoyo psicológico. \\
36--63 & Buscar apoyo profesional de salud mental a la brevedad, evitar situaciones desencadenantes cuando sea posible, contar con red de apoyo cercana y seguir indicaciones de un profesional. \\
\bottomrule
\caption{Recomendaciones por nivel --- BAI}
\label{tab:bai_recomendaciones}
\end{longtable}

\subsection{Declaración de cumplimiento normativo}
\label{sec:cumplimiento_normativo}

El sistema AnxioTest ha sido diseñado para cumplir con las siguientes normativas y estándares:

\begin{itemize}
  \item \textbf{IEEE 830:} Estándar para especificación de requisitos de software.
  \item \textbf{WCAG 2.1 Nivel AA:} Pautas de accesibilidad al contenido web.
  \item \textbf{GDPR:} Reglamento General de Protección de Datos (principios de minimización y transparencia).
  \item \textbf{Ley Orgánica de Protección de Datos Personales del Ecuador:} Marco legal para la protección de datos personales en el contexto ecuatoriano.
\end{itemize}

\subsection{Historial de versiones}
\label{sec:historial_versiones}

\begin{longtable}{p{0.12\linewidth} p{0.18\linewidth} p{0.30\linewidth} p{0.32\linewidth}}
\toprule
\rowcolor{tableheader!20}
\textbf{Versión} & \textbf{Fecha} & \textbf{Autor} & \textbf{Descripción} \\
\midrule
\endfirsthead
\rowcolor{tableheader!20}
\textbf{Versión} & \textbf{Fecha} & \textbf{Autor} & \textbf{Descripción} \\
\midrule
\endhead
\bottomrule
\endfoot
1.0 & 15/05/2025 & Alexander Ñacato & Versión final del documento aprobada. \\
\bottomrule
\caption{Historial de versiones del documento}
\label{tab:historial}
\end{longtable}